\chapter{Catálogo de artefactos}
\label{cap:anexo-artefactos}

Este anexo amplía el inventario resumido de la Tabla~\ref{tab:inventario-runs}. Su
objetivo no es repetir las métricas del Capítulo~6, sino identificar qué fichero permite
auditar cada resultado, dónde se conserva y si forma parte del historial versionado. El
catálogo distingue las ejecuciones oficiales citadas, los análisis auxiliares que sustentan
figuras o comprobaciones y dos ejecuciones de Fase~2 mantenidas únicamente para documentar
la sensibilidad del pipeline a la configuración.

\section{Convenciones de lectura}

Se considera \emph{rastreado} un artefacto devuelto por \texttt{git ls-files} y, por
tanto, recuperable desde el commit de la memoria. Un fichero \emph{local} existe en el
entorno de trabajo, pero está excluido por \texttt{.gitignore}; su existencia no puede
presuponerse en otro clon. La indicación \emph{histórico auditable} significa que se
conservan la configuración y los resultados originales, aunque el código actual no
reproduzca necesariamente la ejecución byte a byte.

Los tipos de fichero tienen funciones diferenciadas:

\begin{itemize}
    \item \texttt{config.json} registra, según el tipo de ejecución, los parámetros y,
    cuando procede, la partición, la semilla y la procedencia de las entradas;
    \item \texttt{metrics.json} o \texttt{validation\_results.json} contiene la salida
    cuantitativa de una ejecución;
    \item \texttt{environment.json} fija el entorno de software observado durante MAIN;
    \item \texttt{scaler.joblib}, \texttt{train\_percentiles.npz} y
    \texttt{feature\_names.json} conservan el estado del preprocesamiento;
    \item \texttt{artifact\_manifest.json} vincula los artefactos confiables mediante
    rutas y resúmenes SHA-256;
    \item los ficheros de predicciones permiten auditar salidas por flujo sin convertir
    una muestra rastreada en sustituto del conjunto completo local.
\end{itemize}

\begingroup
\footnotesize
\setlength{\tabcolsep}{3pt}
\renewcommand{\arraystretch}{1.12}
\begin{longtable}{p{2.7cm}p{6.0cm}p{4.0cm}}
\caption[Catálogo de artefactos de entrenamiento, validación y Fase 2]{Catálogo de
artefactos que respaldan los resultados de la memoria. Las rutas y el estado de
seguimiento proceden del árbol Git de la revisión documentada.}
\label{tab:catalogo-artefactos} \\
\toprule
Elemento & Ruta y artefactos principales & Estado y uso en la memoria \\
\midrule
\endfirsthead
\multicolumn{3}{l}{\small\itshape Continuación de la Tabla~\thetable} \\
\toprule
Elemento & Ruta y artefactos principales & Estado y uso en la memoria \\
\midrule
\endhead
\midrule
\multicolumn{3}{r}{\small\itshape Continúa en la página siguiente} \\
\endfoot
\bottomrule
\endlastfoot

\multicolumn{3}{l}{\textbf{Entrenamiento principal y partición de referencia}} \\
\addlinespace
MAIN QRDQN &
{\raggedright\path{runs/cicids2017/MAIN_qrdqn_cicids2017_canonical_full_random_20260609_193655/}: \texttt{artifact\_manifest.json}, \texttt{config.json}, \texttt{environment.json}, \texttt{feature\_names.json}, \texttt{metrics.json}, \texttt{scaler.joblib}, \texttt{train\_percentiles.npz} y escalares TensorBoard exportados. Modelo: \path{models/MAIN_qrdqn_cicids2017_canonical_full_random_20260609_193655.zip}.\par} &
Oficial y rastreado. Respalda el entrenamiento principal, la evaluación aleatoria fija, el preprocesamiento reutilizado y las curvas de entrenamiento. El evento TensorBoard bruto es local; las exportaciones CSV son la fuente durable. \\
\addlinespace
Referencia de prueba, semilla 42 &
{\raggedright\path{runs/cicids2017/test_partition_reference_seed42.json}.\par} &
Rastreado. Fija conteos y hashes de la partición usada para comprobar que los futuros presupuestos de entrenamiento conservan el mismo conjunto de prueba. \\

\addlinespace
\multicolumn{3}{l}{\textbf{Comprobaciones y análisis de validación}} \\
\addlinespace
VAL A &
{\raggedright\path{runs/validation/VAL_checks_A_20260212_235443/}: \texttt{config.json} y \texttt{validation\_results.json}.\par} &
Oficial citado, rastreado e histórico auditable. Verifica predicción directa; su configuración apunta a un modelo C01 ya retirado del árbol actual. \\
\addlinespace
VAL B &
{\raggedright\path{runs/validation/VAL_checks_B_20260212_235736/}: \texttt{config.json} y \texttt{validation\_results.json}.\par} &
Oficial citado, rastreado e histórico auditable. Registra el control con etiquetas barajadas; el modelo temporal reentrenado no se persistió. \\
\addlinespace
VAL C &
{\raggedright\path{runs/validation/VAL_checks_C_20260213_004847/}: \texttt{config.json} y \texttt{validation\_results.json}.\par} &
Oficial citado, rastreado e histórico auditable. Respalda la partición por día y corresponde a una red proxy, no a los pesos MAIN; el modelo temporal no se persistió. \\
\addlinespace
Bootstrap MAIN &
{\raggedright\path{runs/validation/bootstrap_ci_seed42.json}.\par} &
Rastreado. Contiene $10\,000$ remuestreos estratificados, los conteos de la matriz de confusión y la verificación independiente del modelo MAIN sobre la partición reproducida. \\
\addlinespace
Duplicados &
{\raggedright\path{runs/validation/duplicate_analysis_seed42.json}.\par} &
Rastreado. Cuantifica duplicados exactos y coincidencias entre entrenamiento y prueba para la partición aleatoria de semilla 42. \\

\addlinespace
\multicolumn{3}{l}{\textbf{Línea base Random Forest}} \\
\addlinespace
RF, partición aleatoria &
{\raggedright\path{runs/cicids2017/baseline_random_forest_comparison/rf_cicids2017_canonical_20260628_024735__random_split/}: \texttt{config.json} y \texttt{metrics.json}. Modelo de esta partición: \path{models/rf_cicids2017_canonical_20260628_024735.joblib}.\par} &
Oficial y rastreado. Comparación directa con MAIN sobre la misma partición aleatoria fija; el modelo global conservado corresponde a este barrido. \\
\addlinespace
RF, partición por día &
{\raggedright\path{runs/cicids2017/baseline_random_forest_comparison/rf_cicids2017_canonical_20260628_024735__day_split/}: \texttt{config.json} y \texttt{metrics.json}.\par} &
Oficial y rastreado. Comparación con Check C sobre entrenamiento lunes--miércoles y prueba jueves--viernes. \\
\addlinespace
RF, miércoles retenido &
{\raggedright\path{runs/cicids2017/baseline_random_forest_comparison/rf_cicids2017_canonical_20260628_024735__leave_one_out_wednesday/}: \texttt{config.json} y \texttt{metrics.json}.\par} &
Oficial y rastreado. Resultado exclusivo de RF: no existe todavía un artefacto QRDQN equivalente. \\

\addlinespace
\multicolumn{3}{l}{\textbf{Inferencia offline de Fase 2}} \\
\addlinespace
P2v2 inicial &
{\raggedright\path{runs/phase2/P2v2_pred_20260224_004121/}: \texttt{config.json}, \texttt{metrics.json}, \texttt{diagnostics.json} y \texttt{predictions.csv}.\par} &
Rastreado como contexto histórico. Documenta una tasa de bloqueo extrema bajo una configuración anterior; no respalda el resultado oficial de Fase~2. \\
\addlinespace
P2v2 posterior &
{\raggedright\path{runs/phase2/P2v2_pred_20260408_230318/}: \texttt{config.json}, \texttt{metrics.json} y \texttt{predictions.csv}.\par} &
Rastreado como contexto histórico. Su comportamiento distinto evidencia sensibilidad a la configuración y obliga a citar cada \texttt{RUN\_ID}. \\
\addlinespace
P2v2 MAIN &
{\raggedright\path{runs/phase2/P2v2_pred_20260610_161231_MAIN/}: \texttt{config.json}, \texttt{metrics.json}, \texttt{diagnostics.json} y \texttt{predictions\_head\_10000.csv}. El fichero completo \texttt{predictions.csv.gz} es local.\par} &
Oficial. Configuración, métricas, diagnósticos y muestra de $10\,000$ predicciones rastreados; predicciones completas ignoradas por la regla \texttt{*.gz}. La evidencia se limita al laboratorio cerrado generado por el operador. \\

\end{longtable}
\endgroup

\section{Proveniencia de la entrada de Fase 2}

El fichero \texttt{config.json} de \texttt{P2v2\_pred\_20260610\_161231\_MAIN}
conserva la ruta histórica \path{pcaps/synthetic_real_traffic.csv}. Después de la
ejecución, la entrada se renombró localmente a \path{pcaps/lab_capture_traffic.csv} para
evitar una descripción incorrecta: son paquetes reales capturados en un laboratorio
doméstico cerrado, cuyo tráfico y etiquetas de intención fueron generados por el operador,
no filas de características sintetizadas algorítmicamente. La ruta de la configuración no
se reescribe porque forma parte de la procedencia inmutable de aquella ejecución. El CSV de
entrada y el fichero completo de predicciones permanecen ignorados por su tamaño y por la
sensibilidad potencial de los metadatos; la muestra rastreada sirve para auditoría, no para
reconstruir el conjunto completo.

\section{Elementos excluidos del catálogo oficial}

Las sondas C01--C03, los intentos MAIN rápidos y las demás ejecuciones antiguas de Fase~2 se
conservan bajo \path{runs/archive/} para trazabilidad histórica, pero no respaldan las
cifras oficiales. Del mismo modo, \path{runs/cicids2017/baseline_random_forest_comparison/results_rf.txt}
corresponde a un prototipo supersedido y no se utiliza como evidencia. Los experimentos
QRDQN \textit{leave-one-CSV-out}, la escalera de tamaños de entrenamiento y el estudio
multisemilla no aparecen como ejecuciones porque siguen pendientes de GPU: existen el
código o el protocolo, pero no artefactos comprometidos ni métricas que catalogar.
